\chapter{CONSIDERAÇÕES FINAIS E SUGESTÕES}
\label{capitulo-consideracoes-finais}
	Esta é a parte final do Relatório, na qual o aluno deve apresentar as principais conclusões alcançadas com a atividade desenvolvida.

	Deve expor como a atividade foi importante para sua formação técnica profissional e ressaltar de que maneira os conhecimentos práticos advindos da atividade desenvolvida contribuirão para o conhecimento teórico a ser obtido no curso.

	Com base no trabalho corrente, poderão ser indicados assuntos de relevância significativa para serem abordados em projetos posteriores (sugestões, próximos passos, melhorias, mudanças no encaminhamento das atividades, contribuições do trabalho para a vida acadêmica dos estudantes proponentes da extensão). 
	
    \section{Uso de Inteligência Artificial Generativa}
    \label{capitulo-consideracoes-iag}
    A Portaria CNPq n\textordmasculine  2.664/2026, publicada em 6 de março de 2026, institui a Política de Integridade na Atividade Científica e estabelece as primeiras diretrizes nacionais para o uso de Inteligência Artificial Generativa (IAG) em pesquisas apoiadas pelo CNPq.

    Neste sentido, o uso de Inteligência Artificial Generativa (IAG) pelos alunos será permitida, desde que os alunos informem o seu uso de acordo com os termos da portaria. Em todas as atividades de entrega de trabalhos, os alunos deverão informar os prompts de pesquisa utilizados, qual a IAG utilizada e como ela ajudou no trabalho. Caso seja detectado uso de IAG nas atividades entregues e o aluno não tenha seguido essas regras, o trabalho será considerado inválido e não contará para fins de avaliação. Veja o apêndice \ref{apendice_declaracao_iag} para ver alguns modelos de declaração.

    Outros possíveis declarações para uso de IAG:
    \begin{itemize}
        \item Declaração de uso de IAG para geração/apoio na análise de dados;
        \item Declaração de uso de IAG para geração de código/scripts de pesquisa;
        \item Declaração de uso de IAG para tradução de texto;
        \item Declaração de uso de IAG para geração ou edição de imagens, gráficos e figuras;
        \item Declaração de uso de IAG para transcrição de entrevistas/áudios;
        \item Declaração de uso de IAG para elaboração de resumos/abstracts;
        \item Declaração de uso de IAG na concepção do projeto de pesquisa (brainstorming, formulação de hipóteses);
        \item Declaração de uso de IAG na elaboração de instrumentos de coleta (questionários, roteiros de entrevista);
        \item Declaração de não utilização de IAG (declaração negativa, quando nenhuma ferramenta foi usada);
        \item Declaração de uso de IAG por membros de comitês/pareceristas (uso restrito, vedada inserção de projetos de terceiros);
        \item Declaração de uso de IAG na etapa de submissão (formatação, adequação a normas editoriais).
    \end{itemize}

    Então adapte os modelos sugeridos e faça uma declaração para cada uso de IAG que você.